% !TeX program = xelatex
% Meeting Follow-up Slides #3 -- Gary Wang
% Topic: Amplitude Encoding Ablation (2x2 Factorial)
% Uses NordlingLab169 Beamer theme

\documentclass[aspectratio=169]{beamer}

% Graphics paths must be set BEFORE \usetheme, because the NordlingLab theme
% declares logo images via \pgfdeclareimage at load time.
\usepackage{graphicx}
\graphicspath{
    {../}
    {../../result/multi_model_results/pdf/}
    {../../result/fs_results/pdf/}
    {../../result/fixed_fs_results/pdf/}
    {../Figures/}
}
\makeatletter
\def\input@path{{../}{./}}
\makeatother

\mode<presentation>{
    \usetheme{NordlingLab169}
}

\usepackage[english]{babel}
\let\latinencoding\relax
\usepackage{xltxtra}
\setsansfont{Candara}

\usepackage[absolute,overlay]{textpos}
\usepackage[yyyymmdd]{datetime}
\renewcommand{\dateseparator}{--}

\usepackage{booktabs}
\usepackage{xcolor}
\usepackage{colortbl}
\usepackage{multicol}
\usepackage{tikz}
\usepackage{caption}
\captionsetup[figure]{font=tiny,labelfont=tiny}

\title[Meeting Follow-up \#3]{Progress Update \#3:\\Feature Encoding \& Selection for Text-Only LLMs}
\author[Gary Wang]{Gary Wang}
\institute[NCKU]{Nordling Lab @
    Dept.\ of Mechanical Engineering\\
    National Cheng Kung University\\
    No.\ 1 University rd., Tainan 70101, Taiwan}
\date{\today}

\AtBeginSection[]{
    \begin{frame}<beamer>{Outline}
        \tableofcontents[currentsection]
    \end{frame}
}

%%%%%%%%%%%%%%%%%%%%%%%%%%%%%%%%%%%%%%%%%%%%%%%%%%%%%
\begin{document}

\setbeamertemplate{background}[NLTitle]
\setbeamertemplate{footline}[NLTitle]
\begin{frame}
    \titlepage
\end{frame}

\setbeamertemplate{background}[NLCC]
\setbeamertemplate{footline}[NLCC]
\begin{frame}<beamer>{Outline}
    \tableofcontents
\end{frame}

%%%%%%%%%%%%%%%%%%%%%%%%%%%%%%%%%%%%%%%%%%%%%%%%%%%%%
\section{Motivation: From Code-Executing to Text-Only LLMs}

\begin{frame}
    \frametitle{The New Problem That Offline Extraction Creates}
    Once features are extracted offline, each feature must be
    \emph{rendered textually} in the prompt.
    For every feature we face a design decision:

    \vspace{0.6em}
    \begin{center}
        \textbf{Temporal sequence} \hspace{1em} vs.\hspace{1em}
        \textbf{Aggregated scalar statistic}
    \end{center}

    \vspace{0.8em}
    \textbf{Amplitude is the critical case:}
    \begin{itemize}
        \item MDS-UPDRS 3.4 explicitly references
              \emph{``amplitude decrements near the end / midway /
              after the 1st tap''} --- a temporal property.
        \item Whether scalar summaries \emph{add value on top of}
              the sequence, or merely duplicate it, is not
              obvious a priori --- this is what we test.
    \end{itemize}
\end{frame}

%%%%%%%%%%%%%%%%%%%%%%%%%%%%%%%%%%%%%%%%%%%%%%%%%%%%%
\section{Experiment 1 --- Amplitude Sequence as a Feature Candidate}

\begin{frame}
    \frametitle{Amplitude Sequence as a Feature}
    \footnotesize
    20 \texttt{tap\_amp\_norm} positions + 67 scalars = 87 features ($n{=}53$).
    \begin{columns}[T]
        \begin{column}{0.52\textwidth}
            \centering
            \scriptsize
            \setlength{\tabcolsep}{3pt}
            \begin{tabular}{lcc}
                \toprule
                \textbf{FS method} & \textbf{\#sc} & \textbf{Amp seq?} \\
                \midrule
                Standard LASSO      & 16 & \cellcolor{green!15}4/20 \\
                Group LASSO         & 8  & No \\
                abess best-subset   & 2  & No \\
                Boruta              & 8  & No \\
                RFECV (RF)          & 7  & \cellcolor{green!15}1/20 \\
                Seq.\ Forward (RF)  & 5  & \cellcolor{green!15}3/20 \\
                Stability Sel.      & 5  & \cellcolor{green!15}2/20 \\
                \bottomrule
            \end{tabular}
        \end{column}
        \begin{column}{0.45\textwidth}
            \centering
            \textbf{Top inclusion freq.}\\
            \scriptsize
            \setlength{\tabcolsep}{3pt}
            \begin{tabular}{lc}
                \toprule
                \textbf{Feature} & \textbf{Freq} \\
                \midrule
                amplitude\_min          & 6/7 \\
                period\_quartile\_range & 5/7 \\
                periodEntropy           & 4/7 \\
                amplitude\_stdev        & 4/7 \\
                \cellcolor{green!15}tap\_amp\_norm\_10 & \cellcolor{green!15}3/7 \\
                \cellcolor{green!15}tap\_amp\_norm\_07 & \cellcolor{green!15}2/7 \\
                \cellcolor{green!15}tap\_amp\_norm\_11 & \cellcolor{green!15}2/7 \\
                \bottomrule
            \end{tabular}
        \end{column}
    \end{columns}
    \vspace{0.2em}
    \scriptsize
    \textbf{4/7} methods pick amp-seq positions; late taps (07/10/11) dominate --- matches MDS-UPDRS decrement criterion.\\
    \textbf{$\Rightarrow$} Most methods still \emph{skip} amp-seq entirely --- motivates Exp 2: is amp-seq informative even when FS leaves it out?
\end{frame}

%%%%%%%%%%%%%%%%%%%%%%%%%%%%%%%%%%%%%%%%%%%%%%%%%%%%%
\section{Experiment 2 --- Does the Amplitude Sequence Help?}

\begin{frame}
    \frametitle{Multi-Model Summary (Key Configurations)}
    {\scriptsize
    \setlength{\tabcolsep}{3pt}
    \begin{center}
    \begin{tabular}{llcccc}
        \toprule
        \textbf{Model} & \textbf{Condition} & \textbf{MAE}$\downarrow$ & \textbf{ACC}$\uparrow$ & \textbf{Kappa}$\uparrow$ & \textbf{AAcc}$\uparrow$ \\
        \midrule
        \multicolumn{6}{l}{\textit{DeepSeek-R1 32B FS}} \\
        DS-R1-32b FS & with amp sequence & \cellcolor{green!15}\textbf{0.49} & \cellcolor{green!15}\textbf{0.53} & \cellcolor{green!15}\textbf{0.57} & \cellcolor{green!15}\textbf{0.98} \\
        DS-R1-32b FS & no amp sequence   & 0.55 & 0.49 & 0.55 & 0.96 \\
        \midrule
        \rowcolor{gray!15} Doctor (Chien) & reference & 0.36 & 0.64 & 0.83 & 1.00 \\
        \rowcolor{gray!15} Doctor (Tan)   & reference & 0.40 & 0.60 & 0.69 & 1.00 \\
        \bottomrule
    \end{tabular}
    \end{center}
    }
    \vspace{0.3em}
    \footnotesize
    \begin{itemize}\setlength{\itemsep}{1pt}
        \item Adding the amplitude sequence to the prompt: \textbf{MAE $-0.06$, ACC $+0.04$, Kappa $+0.02$}.
        \item The \emph{with amp sequence} configuration approaches the doctor reference on MAE (0.49 vs.\ Tan 0.40 / Chien 0.36).
    \end{itemize}
    \vspace{0.2em}
    \tiny{amp sequence: \texttt{tap\_amplitudes\_normalized} is rendered in the prompt as a per-tap list.}
\end{frame}

%%%%%%%%%%%%%%%%%%%%%%%%%%%%%%%%%%%%%%%%%%%%%%%%%%%%%
\section{Conclusion}

\begin{frame}
    \frametitle{Summary}
    \begin{itemize}
        \item Extending Chen (2024) to text-only LLMs forced
              us to decide how to encode each feature in text.
        \item Exp 1 --- Amp seq as FS candidate: 4/7 methods
              select late-tap positions (07/10/11), matching
              MDS-UPDRS decrement wording; 3/7 skip it entirely.
        \item Exp 2 --- Adding the amp sequence to the prompt
              gives MAE $-0.06$, ACC $+0.04$; the sequence
              carries independent information.
        \item The \emph{with amp sequence} configuration
              approaches the inter-rater doctor reference on
              MAE.
    \end{itemize}
\end{frame}

\end{document}
